\section{Standard words and their Gram matrices}
\label{sec:standard-grams}

We first fix the word conventions. This is essential because two
standard morphisms can have the same incidence matrix and different
literal images. Throughout, a word is a finite word over $\{a,b\}$,
$\widetilde z$ denotes reversal, and $\eps$ denotes the empty word.
A palindrome satisfies $z=\widetilde z$. Products of matrices follow
the order of letters from left to right.

\subsection{Palindromic closure and standard morphisms}

For a word $z$, let $z^+$ be the shortest palindrome having $z$ as a
prefix. The iterated right palindromization map is defined by
$$
 \operatorname{Pal}(\eps)=\eps,\qquad
 \operatorname{Pal}(dc)=(\operatorname{Pal}(d)c)^+,
 \quad c\in\{a,b\}.
$$
For example, $\operatorname{Pal}(ab)=aba$ and
$\operatorname{Pal}(aba)=abaaba$. We call the values of this map
\emph{central words}. This is the usual binary central-word class;
its relation with Christoffel words is recalled below~\cite[Section~3]{BertheDeLucaReutenauer2008}.
The word $d$ is called a \emph{directive}.

Use the right standard morphisms
\begin{equation}\label{eq:right-standard}
 \phi_a:\ a\mapsto a,\ b\mapsto ba,
 \qquad
 \phi_b:\ a\mapsto ab,\ b\mapsto b,
 \qquad \phi_{dc}=\phi_d\circ\phi_c.
\end{equation}
The corresponding left morphisms are
$\psi_a(a)=a$, $\psi_a(b)=ab$, $\psi_b(a)=ba$, and
$\psi_b(b)=b$. Define, for every word $z$,
$$
 F_c(z)=c\phi_c(z)=\psi_c(z)c,\qquad
 F_{dc}=F_d\circ F_c,\qquad F_\eps(z)=z.
$$
The equality in this definition follows by moving the final $c$
across the image of each letter. Reversing it shows that $F_c$
preserves palindromes. Justin's functional identity~\cite[Lemma~2.1]{Justin2005},
in its right-morphism form~\cite[Section~2]{PerrinReutenauer2023}, identifies
$F_d(\eps)$ with $\operatorname{Pal}(d)$; in the present convention
it gives
\begin{equation}\label{eq:standard-palindrome-image}
 F_d(z)=g_d\phi_d(z),\qquad g_d=\operatorname{Pal}(d),
 \qquad F_d(\operatorname{Pal}(t))=\operatorname{Pal}(dt).
\end{equation}
The first equality also follows directly by induction on $|d|$
once $F_d(\eps)=g_d$ is known. Thus $F_d$ extends the action on
central words to every binary palindrome.

The completed-count formulas below are the incidence form of the
standard central-word construction~\cite[Corollaries~3.1--3.2]{BertheDeLucaReutenauer2008}.
We include the short proof to fix the composition convention.

\begin{lemma}[Literal identities and completed counts]
\label{lem:standard-identities}
For every directive $d$,
\begin{equation}\label{eq:literal-standard-identities}
 \phi_d(ab)=abg_d,\qquad \phi_d(ba)=bag_d.
\end{equation}
Let $L_d$ be the incidence matrix of $\phi_d$, with columns recording
the numbers of $a$'s and $b$'s in the images of $a$ and $b$.
For the \emph{completed count column}
$$
 c(z)=\binom{|z|_a+1}{|z|_b+1},
$$
one has
\begin{equation}\label{eq:completed-count-transport}
 c(F_d(z))=L_dc(z),\qquad
 L_a=\begin{pmatrix}1&1\\0&1\end{pmatrix},\quad
 L_b=\begin{pmatrix}1&0\\1&1\end{pmatrix}.
\end{equation}
In particular, the greatest common divisor of the two completed
counts is unchanged by $F_d$, and every central word has primitive
completed counts.
\end{lemma}

\begin{proof}
The identities hold for the empty directive. Write $u=\phi_d(a)$
and $v=\phi_d(b)$. Appending $a$ changes $(u,v)$ to $(u,vu)$ and
$g_d$ to $g_du$; appending $b$ changes them to $(uv,v)$ and $g_dv$.
If $uv=abg_d$ and $vu=bag_d$, these updates immediately preserve
both identities in \eqref{eq:literal-standard-identities}.
Counting their letters gives
$c(g_d)=L_d(1,1)^{\mathsf T}$. Counting
$F_d(z)=g_d\phi_d(z)$ then proves
\eqref{eq:completed-count-transport}. Each $L_d$ is an integral
unimodular matrix. Such a matrix preserves the ideal generated by
the entries of a column, and therefore their greatest common divisor.
Apply this to $c(\eps)=(1,1)^{\mathsf T}$ for the last assertion.
\end{proof}

Primitivity is only a necessary condition for centrality among
palindromes. We shall use precisely that implication, never its
converse. The maps $F_d$ are injective: each of the two morphisms in
\eqref{eq:right-standard} has a prefix code as its set of letter
images, so is injective; a fixed prefix $g_d$ can then be cancelled.

\subsection{A positive integral Gram model}

The Cohn representation used in this paper is
\begin{equation}\label{eq:cohn-generators}
 \mu(a)=A=\begin{pmatrix}2&1\\1&1\end{pmatrix},\qquad
 \mu(b)=B=\begin{pmatrix}5&2\\2&1\end{pmatrix},\qquad
 P=\begin{pmatrix}2&1\\1&0\end{pmatrix},\quad e=\binom21.
\end{equation}
The matrices $A,B$ are symmetric, positive definite and unimodular.
Consequently $\mu(\widetilde z)=\mu(z)^{\mathsf T}$, and for each
palindrome $z$ the matrix $\mu(z)$ is positive definite.
Indeed, if $z=\widetilde w w$, it is $\mu(w)^{\mathsf T}\mu(w)$;
if $z=\widetilde w c w$, it is
$\mu(w)^{\mathsf T}\mu(c)\mu(w)$.

Define its \emph{Gram label} $M(z)$ and \emph{Gram root} $\rho(z)$ by
\begin{equation}\label{eq:palindromic-gram}
 \mc G(z)=P\mu(z)P
 =\begin{pmatrix}M(z)&\rho(z)\\\rho(z)&H(z)\end{pmatrix}.
\end{equation}
These are positive integers, and
\begin{equation}\label{eq:gram-determinant}
 M(z)H(z)-\rho(z)^2=1.
\end{equation}
In particular, $\rho(z)^2\equiv-1\pmod{M(z)}$.
The label is also a primitive sum of two squares. To see this directly,
put $P_1=\left(\begin{smallmatrix}1&1\\1&0\end{smallmatrix}\right)$
and $P_2=P$. Then $A=P_1^2$, $B=P_2^2$, and the factorizations above
give $\mc G(z)=T^{\mathsf T}T$ with an integral unimodular $T$:
use $T=\mu(w)P$ in the even case and $T=P_c\mu(w)P$ in the odd
case, where $P_a=P_1$ and $P_b=P_2$. The first column of $T$ is
positive and primitive, so its squared Euclidean norm is $M(z)$.

There is also a direct connection with the more familiar trace label.
For $J_*=AP^{-1}=\left(\begin{smallmatrix}1&0\\1&-1\end{smallmatrix}\right)$,
one has
\begin{equation}\label{eq:completed-word-trace}
 \mu(azb)=J_*\mc G(z)P,
 \qquad \tr\mu(azb)=3M(z).
\end{equation}
Both formulas follow by multiplication, using $B=P^2$.
If $z$ is central, $azb$ is a proper lower Christoffel word, and
the classical Cohn correspondence identifies $M(z)$ as a Markov
number~\cite[Theorem~1]{Reutenauer2009}. Here a Markov number is an entry of a positive integral
solution of $x^2+y^2+z^2=3xyz$. For a general palindrome,
\eqref{eq:completed-word-trace} remains an exact identity, but it
does not assert that the label is a Markov number.

Among palindromes with prescribed completed counts, the central word is
distinguished by the label itself. Lagisquet, Pelantov\'a, Tavenas and
Vuillon prove a minimum for the upper-right path value at fixed counts,
with its exact equality cases~\cite[Proposition~21]{LagisquetEtAl2021};
transported to the present Gram normalization, it says that for coprime
completed counts every palindrome satisfies $M(z)\geqslant M(z_0)$,
where $z_0$ is the central word with those counts, and that equality
holds only for $z_0$. We record this because it fixes the direction of
the comparison used below: a competing palindrome with the same counts
can only have a larger label, never a smaller one. The statement is
theirs; nothing in this paper strengthens it, and the noncoprime case is
not used.

\subsection{What the Gram root counts}
\label{subsec:gram-root-counts}

The identity \eqref{eq:gram-determinant} makes $\rho(z)$ a square root
of $-1$ modulo the label. That root also has a word-theoretic meaning,
and the incidence transport of
\eqref{eq:completed-count-transport} is enough to find it. Write
$\widehat d$ for the \emph{reversal complement} of a word $d$, obtained
by reading $d$ backwards and exchanging the two letters.

\begin{lemma}[The reversal complement transposes the incidence matrix]
\label{lem:incidence-transpose}
For every directive $d$, $L_{\widehat d}=L_d^{\mathsf T}$.
\end{lemma}

\begin{proof}
The two matrices of \eqref{eq:completed-count-transport} satisfy
$L_a^{\mathsf T}=L_b$ and $L_b^{\mathsf T}=L_a$, so exchanging the
letters of a directive transposes each factor of its incidence product.
Reading the directive backwards reverses the order of those factors.
Performing both operations therefore transposes the product.
\end{proof}

\begin{proposition}[The doubled directive and its letter counts]
\label{prop:doubled-directive-counts}
Let $d$ be a directive and let
$$
 \binom qs=c(\operatorname{Pal}(d))=L_d\binom11
$$
be its completed counts, so that $\gcd(q,s)=1$. Put
$m=q^2+s^2$. Then the completed counts of the doubled central word
$\operatorname{Pal}(\widehat d\,d)$ are
\begin{equation}\label{eq:doubled-counts}
 c\bigl(\operatorname{Pal}(\widehat d\,d)\bigr)
 =L_d^{\mathsf T}L_d\binom11 ,
\end{equation}
so the completed word $a\operatorname{Pal}(\widehat d\,d)b$ has length
exactly $m$.\footnote{That a Markov number is a sum of two squares was
obtained earlier, by a substitution argument on words resting on Justin's
formula, by L.~Vuillon and A.~De~Luca (personal communication, 2026). The
proof below is independent of theirs and yields the letter counts as well as
the length.} Writing $r_d$ for its number of letters $b$, one has
$0<r_d<m$ and
\begin{equation}\label{eq:doubled-root}
 r_d\,q\equiv s \pmod m,
 \qquad\hbox{hence}\qquad r_d^{\,2}\equiv-1\pmod m .
\end{equation}
The number of letters $a$ is $m-r_d$.
\end{proposition}

\begin{proof}
By \eqref{eq:standard-palindrome-image},
$\operatorname{Pal}(\widehat d\,d)=F_{\widehat d}(\operatorname{Pal}(d))$,
so \eqref{eq:completed-count-transport} and
\cref{lem:incidence-transpose} give \eqref{eq:doubled-counts}. Write
$L_d=\left(\begin{smallmatrix}\alpha&\beta\\\gamma&\delta\end{smallmatrix}\right)$,
so that $q=\alpha+\beta$ and $s=\gamma+\delta$, and
$\alpha\delta-\beta\gamma=1$. Expanding \eqref{eq:doubled-counts},
$$
 m-r_d=\alpha q+\gamma s,\qquad r_d=\beta q+\delta s,
$$
whose sum is $(\alpha+\beta)^2+(\gamma+\delta)^2=q^2+s^2$, which is the
length assertion.

Both $\alpha$ and $\delta$ are at least one, since $L_d$ has nonnegative
entries and determinant one, and $\alpha=0$ would force
$-\beta\gamma=1$. Hence $r_d\geqslant\delta s\geqslant1$ and
$m-r_d\geqslant\alpha q\geqslant1$, which is the strict range.

Finally $\gcd(q,m)=\gcd(q,s^2)=1$, so $q$ is invertible modulo $m$, and
$q^2\equiv-s^2$. Therefore
$$
 r_d\,q=\beta q^2+\delta sq
 \equiv-\beta s^2+\delta sq
 =s\bigl(\delta(\alpha+\beta)-\beta(\gamma+\delta)\bigr)
 =s(\alpha\delta-\beta\gamma)=s \pmod m .
$$
Squaring and cancelling the unit $q^2$ gives $r_d^{\,2}\equiv-1$.
\end{proof}

The two letter counts of that word are not merely arithmetic data: they are
also the two lengths into which it splits. Recall that a proper Christoffel
word has a unique factorization into two Christoffel words, its
\emph{standard factorization}, and that the lengths of the two factors are the
column sums of the incidence matrix of its
directive~\cite[Corollaries~3.1--3.2]{BertheDeLucaReutenauer2008}.

\begin{corollary}[The cut and the counts are the same pair]
\label{cor:cut-at-letter-count}
Let $d$ be a directive, let
$\Omega=a\operatorname{Pal}(\widehat d\,d)b$, and let
$\Omega=\Omega_1\Omega_2$ be its standard factorization. Then
\begin{equation}\label{eq:cut-at-letter-count}
 \abs{\Omega_1}=\abs{\Omega}_a=m-r_d,
 \qquad
 \abs{\Omega_2}=\abs{\Omega}_b=r_d .
\end{equation}
That is, $\Omega$ is cut after exactly as many letters as it contains
letters $a$.
\end{corollary}

\begin{proof}
The directive of $\Omega$ is $\widehat d\,d$, so by
\cref{lem:incidence-transpose} its incidence matrix is
$L_{\widehat d\,d}=L_{\widehat d}L_d=L_d^{\mathsf T}L_d$. This is a Gram
matrix, hence symmetric, and therefore its column sums and its row sums are
the same ordered pair:
$$
 \mathbf 1^{\mathsf T}L_d^{\mathsf T}L_d
 =\bigl(L_d^{\mathsf T}L_d\mathbf 1\bigr)^{\mathsf T},
 \qquad \mathbf 1=(1,1)^{\mathsf T}.
$$
The left side is the pair of factor lengths of the standard factorization,
and the right side is, by \eqref{eq:doubled-counts}, the pair of completed
counts of $\Omega$. The identification of the individual entries is the
order of \eqref{eq:doubled-counts}, and $\abs{\Omega}_a=m-r_d$ was proved in
\cref{prop:doubled-directive-counts}.
\end{proof}

Symmetry of a Gram matrix is the whole content: the cut point and the letter
count are one pair of integers read along the two axes of a single matrix.
Exchanging the two letters of $d$ replaces that matrix by the Gram matrix of
the rows of $L_d$, which has the same total mass and the two entries of
\eqref{eq:cut-at-letter-count} interchanged, so the cut is reflected in the
middle of the word.

The Gram model sees the same root. Recall the factorization
$\mc G(z)=T^{\mathsf T}T$ constructed above, with $T=\mu(w)P$ in the even
case and $T=P_c\mu(w)P$ in the odd case.

\begin{corollary}[The Gram root as a letter count]
\label{cor:gram-root-letter-count}
Let $z$ be a palindrome, let $(q,s)$ be the first column of $T$, and put
$\eps_z=\det T$, which is $+1$ when $\abs z$ is odd and $-1$ when
$\abs z$ is even. Then $M(z)=q^2+s^2$ with $\gcd(q,s)=1$, and
$$
 \rho(z)\,q\equiv\eps_z\,s \pmod{M(z)} .
$$
Let $d$ be the unique directive with completed counts $(q,s)$. Then
$\rho(z)=r_d$ when $\eps_z=+1$, and $\rho(z)=M(z)-r_d$ when
$\eps_z=-1$.
\end{corollary}

\begin{proof}
Let $(\beta,\delta)$ be the second column of $T$. Then $M(z)=q^2+s^2$
and $\rho(z)=\beta q+\delta s$ are the two entries of the first row of
$T^{\mathsf T}T$, and $\det T=q\delta-s\beta=\eps_z$. Since
$\det\mu(w)=1$, $\det P=\det P_1=-1$, the two cases of the construction
give $\eps_z=-1$ and $\eps_z=+1$ respectively. As in the proposition,
$q^2\equiv-s^2$ modulo $M(z)$, so
$$
 \rho(z)\,q=\beta q^2+\delta sq\equiv s(q\delta-s\beta)=\eps_z s .
$$
The first column of $T$ is positive and primitive, so $(q,s)$ is the
completed-count column of exactly one central word, namely
$\operatorname{Pal}(d)$ for a unique directive $d$; this is the
inverse of the unimodular transport
\eqref{eq:completed-count-transport}. Both $\rho(z)$ and $r_d$ lie
strictly between $0$ and $M(z)$, and by
\eqref{eq:doubled-root} they satisfy the same congruence up to the sign
$\eps_z$. A residue class modulo $M(z)$ meets that range once, so they
agree when $\eps_z=+1$ and are complementary when $\eps_z=-1$.
\end{proof}

When $z$ is central, $azb$ is a proper lower Christoffel word and $M(z)$
is a Markov number, so $\rho(z)$ is then the classical root of $-1$
attached to the occurrence, and $q^2+s^2$ is its primitive two-square
representation. \Cref{prop:doubled-directive-counts} says that this root
is not an abstract residue: it is the number of letters $b$ in one
explicit Christoffel word of length $M(z)$, the one directed by the
doubled directive $\widehat d\,d$. The statement itself needs no
centrality, and its proof uses only \eqref{eq:completed-count-transport}
and a determinant. \Cref{cor:cut-at-letter-count} adds that the same integer
also measures where that word splits: the standard factorization cuts after
$M(z)-\rho(z)$ letters when $\eps_z=+1$, and after $\rho(z)$ letters when
$\eps_z=-1$. A root of $-1$ modulo a Markov number is thus visible twice in a
single word, once as a count and once as a position.

\subsection{The Thue--Morse coding}
\label{subsec:thue-morse}

The label and the root have a second reading, in the coding of Markov numbers by Reutenauer
and Vuillon~\cite{ReutenauerVuillon2017}, and a search for equal lengths among all palindromes of
that coding finds the pairs of \cref{sec:midpoint}. Let $\theta$ be
the Thue--Morse morphism, $\theta(a)=ab$ and $\theta(b)=ba$, and recall the matrix
$P_1=\left(\begin{smallmatrix}1&1\\1&0\end{smallmatrix}\right)$ of the Gram factorization, with
$P_1^2=A$.

\begin{proposition}[The Thue--Morse coding]
\label{prop:thue-morse}
For every word $z$,
\begin{equation}\label{eq:thue-morse-incidence}
 L_{\theta(F_a(z))}=P_1\,\mu(z)\,P_1 .
\end{equation}
Consequently, for every palindrome $z$ the Christoffel word
$\Omega_z=a\operatorname{Pal}(\theta(F_a(z)))b$ has length $M(z)$ and contains exactly
$\rho(z)$ letters $b$. Moreover, $F_a$ is a bijection from the set of palindromes onto the
set of palindromes that begin with $a$ and do not contain the factor $bb$.
\end{proposition}

\begin{proof}
The standard morphism of a concatenation of directives is the composition of their
standard morphisms, so $L_{uv}=L_uL_v$; since $\theta$ is a morphism, it suffices to compare
the images of the generators. On the word side, $L_{\theta(a)}=L_aL_b=A$ and
$$
 L_{\theta(\phi_a(b))}=L_{baab}=(L_bL_a)(L_aL_b)
 =\begin{pmatrix}1&1\\1&2\end{pmatrix}\begin{pmatrix}2&1\\1&1\end{pmatrix}
 =\begin{pmatrix}3&2\\4&3\end{pmatrix}.
$$
On the matrix side, $P_1^{-1}AP_1=A$ because $A=P_1^2$, and
$$
 P_1^{-1}BP_1=\begin{pmatrix}0&1\\1&-1\end{pmatrix}\begin{pmatrix}7&5\\3&2\end{pmatrix}
 =\begin{pmatrix}3&2\\4&3\end{pmatrix}.
$$
Since $\theta(\phi_a(a))=\theta(a)$, this proves $L_{\theta(\phi_a(c))}=P_1^{-1}\mu(c)P_1$ for
both letters $c$. The product over the letters of $z$ telescopes to
$L_{\theta(\phi_a(z))}=P_1^{-1}\mu(z)P_1$, and therefore
$L_{\theta(F_a(z))}=L_{\theta(a)}L_{\theta(\phi_a(z))}=P_1^2P_1^{-1}\mu(z)P_1$, which is
\eqref{eq:thue-morse-incidence}.

Every directive $D$ satisfies $c(\operatorname{Pal}(D))=L_D(1,1)^{\mathsf T}$, by
\eqref{eq:standard-palindrome-image} and \eqref{eq:completed-count-transport} applied to the
empty word. Take $D=\theta(F_a(z))$ and write $\mu(z)e=(x,y)^{\mathsf T}$. Since
$P_1(1,1)^{\mathsf T}=e$, the identity \eqref{eq:thue-morse-incidence} gives
$$
 c(\operatorname{Pal}(D))=P_1\mu(z)e=\binom{x+y}{x}.
$$
For a palindrome $z$ one has $M(z)=e^{\mathsf T}\mu(z)e=2x+y$, because the first column of
$P$ is $e$, and $\rho(z)=(1,0)\,\mu(z)e=x$, because the second row of $P$ is $(1,0)$. The
completed counts of $\operatorname{Pal}(D)$ are the letter counts of $\Omega_z$, which are
therefore $M(z)-\rho(z)$ and $\rho(z)$. The word $\Omega_z$ is a Christoffel word because
$\operatorname{Pal}(D)$ is central.

For the last assertion, recall $F_a(z)=a\phi_a(z)=\psi_a(z)a$. The reversal of
$\phi_a(z)$ is $\psi_a(\widetilde z)$, so the reversal of $F_a(z)$ is $F_a(\widetilde z)$. If
$z$ is a palindrome, $F_a(z)$ is thus a palindrome; it begins with $a$, and it has no factor
$bb$ because every $b$ of $\phi_a(z)$ is followed by an $a$. Conversely, let $u=au'$ be a
palindrome without the factor $bb$. The word $u'$ is empty or ends with $a$, so each of its
letters $b$ is followed by the letter $a$, and $u'$ factors uniquely over the code
$\{a,ba\}$: $u'=\phi_a(z)$ and $u=F_a(z)$. Finally $F_a(z)=u=F_a(\widetilde z)$, and the
injectivity of $F_a$ gives $z=\widetilde z$.
\end{proof}

For a central word $z=\operatorname{Pal}(v)$ one has $F_a(z)=\operatorname{Pal}(av)$ by
\eqref{eq:standard-palindrome-image}. Then $\Omega_z$ is the word
$a\operatorname{Pal}(\theta(\operatorname{Pal}(av)))b$ of Reutenauer and Vuillon, and the length
statement is their theorem that $2+\abs{\operatorname{Pal}(\theta(\operatorname{Pal}(av)))}$
is the Markov number of the occurrence~\cite[Theorem~1]{ReutenauerVuillon2017}. The
proposition adds the root. By \eqref{eq:completed-word-trace} the first row of the Cohn
matrix $\mu(azb)=J_*\mc G(z)P$ is $(M,\rho)P=(2M+\rho,\,M)$, so the number of letters $b$ of
their word is $\mu(azb)_{11}-2M$, for every occurrence and with no choice of sign. For an
arbitrary palindrome, $F_a(z)$ is an arbitrary palindrome that begins with $a$ and avoids
$bb$, and a coincidence of the lengths
$\abs{\operatorname{Pal}(\theta(F_a(z)))}$ is exactly an equality of Gram labels. The
pair of the introduction is the first one: with $F_a(baab)=abaaaba$ and $F_a(abba)=aababaa$,
the directives $\theta(abaaaba)$ and $\theta(aababaa)$ give central words whose completed counts
are $(663,467)$ and $(693,437)$, and both completed words have length $1130$.

The Thue--Morse directive is the doubled directive of
\cref{prop:doubled-directive-counts}, which explains the sign in
\cref{cor:gram-root-letter-count}.

\begin{corollary}[The Thue--Morse directive is a doubled directive]
\label{cor:thue-morse-doubled}
Let $z$ be a palindrome and let $d$ be the directive of \cref{cor:gram-root-letter-count},
whose completed counts are the first column of $T$. Then $\theta(F_a(z))=\widehat d\,d$ when
$\abs z$ is odd, and $\theta(F_a(z))$ is obtained from $\widehat d\,d$ by exchanging the two
letters when $\abs z$ is even. In particular the sign $\eps_z$ of
\cref{cor:gram-root-letter-count} is the exchange of letters, and the Gram root is the number
of letters $b$ of $\Omega_z$ in both parities.
\end{corollary}

\begin{proof}
Every matrix of $\SL_2(\Z)$ with nonnegative entries is a product of $L_a$ and $L_b$ in
exactly one way, and the column $L_{d'}(1,1)^{\mathsf T}$ determines the directive $d'$, as in
the proof of \cref{cor:gram-root-letter-count}. In the odd case, $z=\widetilde wcw$, the
matrix $L=P_c\mu(w)P_1$ has determinant $(-1)\cdot(-1)=1$ and nonnegative entries, and
$L(1,1)^{\mathsf T}=P_c\mu(w)e=Te_1$. Hence $L=L_d$, and by \cref{lem:incidence-transpose}
$$
 L_{\widehat d\,d}=L_d^{\mathsf T}L_d=P_1\mu(\widetilde w)P_c^2\mu(w)P_1=P_1\mu(z)P_1,
$$
because $P_c$ is symmetric and $P_c^2=\mu(c)$. In the even case, $z=\widetilde ww$, the
matrix $L=\mu(w)L_a$ has determinant one and nonnegative entries, and
$L(1,1)^{\mathsf T}=\mu(w)e=Te_1$; hence $L=L_d$ and
$L_{\widehat d\,d}=L_b\mu(\widetilde w)\mu(w)L_a=L_b\mu(z)L_a$. With
$S=\left(\begin{smallmatrix}0&1\\1&0\end{smallmatrix}\right)$ one has $P_1=L_aS=SL_b$, so
$P_1\mu(z)P_1=S\,L_{\widehat d\,d}\,S$, and conjugation by $S$ exchanges $L_a$ and $L_b$, that
is, the two letters of a directive. In both cases \eqref{eq:thue-morse-incidence} and the
uniqueness of the factorization prove the claim. Exchanging the letters of a directive
exchanges the two letter counts of its central word, which gives the statement about
$\eps_z$ by comparison with \cref{cor:gram-root-letter-count}.
\end{proof}

The coding turns the question of this article into a question about the lengths of central
words. The pairs classified in \cref{sec:midpoint} are coincidences of these lengths, and
\cref{sec:outlook} asks whether they are the only ones.

\subsection{The reflection attached to an ancestor}

Fix a directive $d$ and abbreviate
\begin{equation}\label{eq:ancestor-data}
 G=\mu(g_d),\quad \nu_d=\mu\circ\phi_d,\quad
 M_0=e^{\mathsf T}Ge,\quad \rho_0=(Ge)_1,\quad r_0=\rho_0/M_0.
\end{equation}
The entries of $Ge$ are positive, so $0<r_0<1/2$.
Equation~\eqref{eq:standard-palindrome-image} becomes
\begin{equation}\label{eq:image-gram}
 \mc G(F_d(z))=PG\nu_d(z)P.
\end{equation}
In particular, $G\nu_d(z)$ is symmetric for a palindrome $z$,
although $\nu_d(z)$ itself need not be symmetric.

The ancestor determines the rational matrix
\begin{equation}\label{eq:ancestor-reflection}
 R_d=\frac{2ee^{\mathsf T}G}{M_0}-I.
\end{equation}
It fixes $e$, squares to $I$, and satisfies
$R_d^{\mathsf T}G=GR_d$ and $R_d^{\mathsf T}GR_d=G$.
Thus it is an orthogonal reflection for the inner product defined by
$G$. For example, if $G=T^{\mathsf T}T$, then $TR_dT^{-1}$ is the
ordinary Euclidean reflection in the line spanned by $Te$.
In the Gram coordinates its form is particularly simple:
\begin{equation}\label{eq:gram-reflection}
 R_d=PJ_0P^{-1},\qquad
 J_0=\begin{pmatrix}1&2r_0\\0&-1\end{pmatrix}.
\end{equation}
The next section determines exactly when this reflection carries a
finite Cohn word column to another such column.
